% Generated by scripts/assemble.py; corrections belong in source/corrections.json.
\examhead{16}{中国科学院大学 2024 年数学分析试卷(回忆版)}{21}
% p022-b01; source PDF page 22
\par\Needspace{5\baselineskip}
\blockanchor{p022-b01}
\majorquestion{一}(15 分) 解答题.\par
% p022-b02; source PDF page 22
\blockanchor{p022-b02}
(1) 证明: $\left|e - \left(1 + \frac{1}{n}\right)^n\right| < \frac{3}{n}, n \in \mathbb{N}^+$.\par
% p022-b03; source PDF page 22
\blockanchor{p022-b03}
(2) 设非负数列 $\{a_{n}\}$ 满足 $a_{n + m} \leq a_{n} + a_{m}$, 证明: $\lim_{n \to \infty} \frac{a_{n}}{n} = \inf \left\{\frac{a_{n}}{n}\right\}$.\par
% p022-b04; source PDF page 22
\blockanchor{p022-b04}
(3) 求极限 $\lim_{x\to +\infty}\frac{\left(\int_0^xe^{t^2}\mathrm{d}t\right)^2}{\int_0^xe^{2t^2}\mathrm{d}t}.$\par
% p022-b05; source PDF page 22
\par\Needspace{5\baselineskip}
\blockanchor{p022-b05}
\majorquestion{二}（15分）设 $f(x)$ 在 $(a,b)$ 上连续，证明以下条件等价：\par
% p022-b06; source PDF page 22
\blockanchor{p022-b06}
(1) $f(x)$ 在 $(a, b)$ 上一致连续.\par
% p022-b07; source PDF page 22
\blockanchor{p022-b07}
(2) $f(x)$ 在端点 $a$ 和 $b$ 处极限存在.\par
% p022-b08; source PDF page 22
\blockanchor{p022-b08}
(3) $f(x)$ 可延拓成 $[a, b]$ 上的连续函数.\par
% p022-b09; source PDF page 22
\par\Needspace{5\baselineskip}
\par\Needspace{6\baselineskip}
\blockanchor{p022-b09}
\majorquestion{三}(15 分) 设 $f(x)$ 在 $[a, b]$ 上二阶可导, 且 $f'(a) = f'(b) = 0$, 证明: 存在 $\xi \in (a, b)$, 使得\par
\nopagebreak[4]
% p022-b10; source PDF page 22
\blockanchor{p022-b10}
\displayformula{| f ^ {\prime \prime} (\xi) | \geq \frac {4}{(b - a) ^ {2}} | f (b) - f (a) |.}
% p022-b11; source PDF page 22
\par\Needspace{5\baselineskip}
\blockanchor{p022-b11}
\majorquestion{四}(15 分) 设 $f(x)$ 在 $[0, 1]$ 上连续, 求极限 $\lim_{n \to \infty} \int_0^1 f(\sqrt[n]{x}) \, \mathrm{d}x$.\par
% p022-b12; source PDF page 22
\par\Needspace{5\baselineskip}
\par\Needspace{6\baselineskip}
\blockanchor{p022-b12}
\majorquestion{五}（15分）设 $f(x,y)$ 为二阶连续可微函数，且 $\frac{\partial^2f}{\partial x^2} +\frac{\partial^2f}{\partial y^2} = x^2 y^2$ ，求重积分\par
\nopagebreak[4]
% p022-b13; source PDF page 22
\blockanchor{p022-b13}
\displayformula{\iint_ {x ^ {2} + y ^ {2} \leq 1} \left(\frac {x}{\sqrt {x ^ {2} + y ^ {2}}} \frac {\partial f}{\partial x} + \frac {y}{\sqrt {x ^ {2} + y ^ {2}}} \frac {\partial f}{\partial y}\right) d x d y.}
% p022-b14; source PDF page 22
\par\Needspace{5\baselineskip}
\blockanchor{p022-b14}
\majorquestion{六}(15 分) 估计 $\ln 2$ 的近似值, 精确到 0.0001.\par
% p022-b15; source PDF page 22
\par\Needspace{5\baselineskip}
\par\Needspace{6\baselineskip}
\blockanchor{p022-b15}
\majorquestion{七}(15 分) 设 $f$ 为 $\mathbb{R}$ 上的凸函数, $\varphi$ 为闭区间 $E$ 上的连续函数, 证明:\par
\nopagebreak[4]
% p022-b16; source PDF page 22
\blockanchor{p022-b16}
\displayformula{f \left(\frac {1}{\mu (E)} \int_ {E} \varphi (x) \mathrm{d} x\right) \leq \frac {1}{\mu (E)} \int_ {E} f \circ \varphi (x) \mathrm{d} x.}
% p022-b17; source PDF page 22
\blockanchor{p022-b17}
其中  $\mu(E)$  为 E 的长度.(可能有误)\par
% p022-b18; source PDF page 22
\par\Needspace{5\baselineskip}
\par\Needspace{6\baselineskip}
\blockanchor{p022-b18}
\majorquestion{八}（15分）设 $I = \iiint_{V} (x + y - z + 10) \, \mathrm{d}V$，其中 $V: x^2 + y^2 + z^2 \leq 3$，证明：\par
\nopagebreak[4]
% p022-b19; source PDF page 22
\blockanchor{p022-b19}
\displayformula{28 \sqrt {3} \pi \leq I \leq 52 \sqrt {3} \pi . (\text { 可   能   有   误 })}
% p022-b20; source PDF page 22
\par\Needspace{5\baselineskip}
\par\Needspace{6\baselineskip}
\blockanchor{p022-b20}
\majorquestion{九}（15分）计算曲线积分\par
\nopagebreak[4]
% p022-b21; source PDF page 22
\blockanchor{p022-b21}
\displayformula{I = \oint_ {L} (y ^ {2} - z ^ {2}) \mathrm{d} x + (z ^ {2} - x ^ {2}) \mathrm{d} y + (x ^ {2} - y ^ {2}) \mathrm{d} z.}
% p022-b22; source PDF page 22
\par\Needspace{6\baselineskip}
\blockanchor{p022-b22}
其中 $L$ 是 $[0, a] \times [0, a] \times [0, a]$ 的表面与平面 $x + y + z = \frac{3}{2} a$ 的交线，从上往下看取逆时针方向。
\majorquestion{十}（15分）证明：光滑曲线 $y = f(x) (f(x) > 0)$ 绕 $x$ 轴旋转一周所得旋转曲面面积为\par
\nopagebreak[4]
% p022-b23; source PDF page 22
\blockanchor{p022-b23}
\displayformula{S = 2 \pi \int_ {a} ^ {b} f (x) \sqrt {1 + (f ^ {\prime} (x)) ^ {2}} \mathrm{d} x.}
